\chapter{State of the art}
\label{chap:state-of-the-art}

This chapter reviews the four bodies of work this thesis draws on: benchmarks of language-model
inference on \glspl{sbc}, post-training quantisation and its evaluation, grammar-constrained
decoding, and spoken-language understanding. Each section first sorts its studies into the kinds
that exist, then compares them to one another in a table along fixed axes, and closes by reading
that table: what the studies agree on, where they differ, and what none of them measures.
Section~\ref{sec:research-gaps} combines the four readings into three numbered gaps and maps each
onto the contribution of Chapter~\ref{chap:introduction} that closes it. The studies are compared
to each other in the tables and to this thesis only in that final section. The foundations the
tables assume, from the K-quant formats to token masking under a grammar, are those of
Chapter~\ref{chap:background}.

\section{Edge language-model inference and single-board-computer benchmarking}
\label{sec:soa-edge}
Running a language model without a network connection or a data-centre accelerator is now an
active subject of measurement, and the studies that measure it are of two kinds. The first kind
benchmarks the hardware: it times many models on a board and reports throughput and resource use
without asking what the tokens are for. A recent benchmark of this kind evaluates twenty-five
quantised language models on three \glspl{sbc} under two inference runtimes, Ollama and
Llamafile~\cite{sbc2025}. It finds that such boards reliably support models of up to roughly
1.5~billion parameters, with larger models dropping to a few tokens per second. On the
one board where both runtimes were compared, Llamafile reached up to four times the throughput of
Ollama~\cite{sbc2025}. That result sets the shape of the problem this thesis inherits. On
\gls{cpu}-only edge hardware throughput is the scarce resource, and the models that fit within it
are small ones; the inference runtime is a lever on it of the same order as the choice of model.

The second kind benchmarks a runtime's formats on one model and scores the output: a unified
evaluation of \texttt{llama.cpp}'s quantisation formats on Llama-3.1-8B-Instruct reports downstream
benchmark accuracy and perplexity next to \gls{cpu} throughput~\cite{kurt2026}. It is the study
nearest in instrumentation to this one, since it uses the same runtime and the same format family,
and Section~\ref{sec:soa-quantisation} returns to it. Table~\ref{tab:soa-edge} sets the two kinds
beside each other and beside this work.

\begin{table}[htbp]
\centering
\footnotesize
\caption[Studies of language-model inference without an accelerator]{Studies of language-model
inference on \gls{cpu}-only hardware, by what they run and what they score. \emph{Task scored}:
whether the study measures any property of the output beyond speed and resource use. A dash marks
a property this review did not establish from the source. The last row is this work.}
\label{tab:soa-edge}
\begin{tabularx}{\textwidth}{@{}>{\raggedright\arraybackslash}p{2.3cm} >{\raggedright\arraybackslash}p{2.4cm} >{\raggedright\arraybackslash}p{2.5cm} >{\raggedright\arraybackslash}X >{\raggedright\arraybackslash}X@{}}
\toprule
Study & Models & Hardware & Runtime and formats & Metrics and task scored \\
\midrule
Nguyen and Nguyen~\cite{sbc2025} & 25 quantised models, up to a few billion parameters & Three \glspl{sbc} & Ollama and Llamafile; formats as distributed & Throughput, resource use; no task \\
Kurt~\cite{kurt2026} & Llama-3.1-8B-Instruct & \gls{cpu}; --- & \texttt{llama.cpp}; K-quant and legacy \gls{gguf} formats against fp16 & Perplexity, generic benchmark accuracy, \gls{cpu} throughput; generic tasks \\
\midrule
This work & Three fine-tuned models of 0.36--1.2~B, plus a control & Raspberry~Pi~5, three \gls{cpu} cores & \texttt{llama.cpp}; Q8\_0 and Q4\_K\_M against fp16 & \gls{em}, intent and slot F1, schema validity, false-command rate, throughput, latency, memory; command parsing under a grammar \\
\bottomrule
\end{tabularx}
\end{table}

\paragraph{Critical comparison.} Table~\ref{tab:soa-edge} separates the two studies on three
axes. On breadth against depth, the hardware benchmark covers twenty-five models on three boards
and scores none of their output, while the format evaluation covers one model on one class of
processor and scores it on several suites; neither covers several models scored on one task. On
the runtime, the two agree that the inference engine is a variable of the same order as the model,
one by comparing engines and the other by comparing one engine's formats. On scale, both stop
short of the sub-billion range: the hardware benchmark reports that boards serve models of up to
about 1.5~billion parameters without saying what accuracy those models retain, and the format
evaluation studies an 8-billion-parameter model that no \gls{sbc} serves at interactive speed. A
configuration that clears a throughput floor and one that clears it while producing wrong output
are therefore indistinguishable in the first study's tables and absent from the second's. This
thesis narrows the scope to one task, one device, and four fine-tuned models
(Qwen2.5-0.5B~\cite{qwen25}, SmolLM2-360M~\cite{smollm2}, Llama-3.2-1B~\cite{llama32}, and
H2O-Danube3-500M~\cite{danube3}), and in exchange scores what neither study does. Three of those
models are carried through to quantised artefacts and timed against the 20~tok/s throughput floor
of Chapter~\ref{chap:introduction}. The fourth, a parameter-matched control, is reported at fp16
only, because \texttt{llama.cpp} segments its prompt differently from the tokeniser it was trained
under (Chapter~\ref{chap:method}). The accuracy cost of meeting a throughput budget on an
\gls{sbc} is thus treated as a quantity to be measured, not assumed from a benchmark that reports
throughput alone. The runtime that makes the measurement possible is \texttt{llama.cpp}~\cite{llamacpp},
described in Section~\ref{sec:bg-edge}; Section~\ref{sec:soa-quantisation} takes up what its
formats cost in accuracy.

\section{Quantisation}
\label{sec:soa-quantisation}
Post-training quantisation is the established route to reducing a trained model's memory
footprint~\cite{gptq}, and on-device deployment is one of its stated motivations~\cite{awq}. The
literature divides into the methods, which propose a way of assigning bits, and the evaluations,
which measure what an assignment costs. This thesis deploys two of \texttt{llama.cpp}'s
\gls{gguf} formats, Q8\_0 and Q4\_K\_M, whose construction Section~\ref{sec:bg-quant} describes.
What distinguishes them from the methods the literature studies is how bits are assigned: the
K-quant recipe allocates them by tensor role, position and shape, not by any measurement of error.
The method literature measures instead: it finds that protecting the roughly one percent of weight
channels that activation statistics mark as salient greatly reduces quantisation
error~\cite{awq}.

The runtime's documentation states that the accuracy cost of its formats is usually measured as
a perplexity or a Kullback--Leibler divergence against the unquantised model~\cite{llamacpp},
not as a change in accuracy on a downstream task. A perplexity number averaged over held-out text
says how much less confident the quantised model is about the next token in general. It does not
say whether the intent assigned to a spoken command, or the slot values filled in, still match the
ground truth once the weights have been rounded to lower precision. This thesis needs the
second number, because its model's output is acted on by a swarm controller after automatic
validation, not read by a person. The papers introducing post-training quantisation methods report
their cost mainly in generic terms: SpQR reports a relative perplexity loss of under one
percent~\cite{spqr}, and GPTQ reports perplexity together with zero-shot accuracy on generic
benchmarks such as LAMBADA, ARC and PIQA~\cite{gptq}.

Evaluation studies have since measured downstream accuracy under quantisation at
scale~\cite{kurtic2025}, and for \texttt{llama.cpp}'s own formats specifically~\cite{kurt2026}. Their
models are mostly of eight billion parameters and above, scored on general-purpose reasoning,
knowledge and coding suites with the decoder unconstrained; the smallest model either study
evaluates has 1.5~billion parameters~\cite{kurtic2025}. Small models have also been studied
directly: a benchmark of quantisation at sub-billion scale, which includes Qwen2.5-0.5B, finds that
small models differ from large ones in their sensitivity to quantisation, so that methods tuned on
large models transfer poorly~\cite{slmquant}. Table~\ref{tab:soa-quantisation} places the methods
and the evaluations on one set of axes.

\begin{table}[htbp]
\centering
\footnotesize
\caption[Quantisation methods and evaluations]{Post-training quantisation methods and
evaluations, by what they quantise and how they report the cost. \emph{Kind}: a method proposes
a bit assignment; an evaluation measures assignments made by others; the runtime is the
implementation deployed here. A dash marks a property this review did not establish from the
source. The last row is this work.}
\label{tab:soa-quantisation}
\begin{tabularx}{\textwidth}{@{}>{\raggedright\arraybackslash}p{2.2cm} >{\raggedright\arraybackslash}p{1.7cm} >{\raggedright\arraybackslash}p{2.2cm} >{\raggedright\arraybackslash}X >{\raggedright\arraybackslash}X >{\raggedright\arraybackslash}p{1.6cm}@{}}
\toprule
Study & Kind & Models and scale & Bit assignment & Cost reported as & Grammar \\
\midrule
GPTQ~\cite{gptq} & Method & --- & One-shot weight quantisation to 3--4 bits with error compensation from calibration data & Perplexity; zero-shot accuracy on generic benchmarks & No \\
AWQ~\cite{awq} & Method & --- & Weight-only; salient channels protected by activation statistics & --- & No \\
SpQR~\cite{spqr} & Method & --- & Outlier weights isolated at higher precision & Relative perplexity loss under one percent & No \\
\texttt{llama.cpp} K-quants~\cite{llamacpp} & Runtime & Any \gls{gguf} model & Bits by tensor role, position and shape; no calibration & Perplexity; Kullback--Leibler divergence & Optional \\
Kurtic et al.~\cite{kurtic2025} & Evaluation & Llama-3.1 family; smallest 1.5~B & FP8, INT8, INT4 & Accuracy on academic benchmarks and real-world tasks & No \\
Kurt~\cite{kurt2026} & Evaluation & Llama-3.1-8B-Instruct & \texttt{llama.cpp} formats from Q3\_K to Q8\_0 & Perplexity; generic benchmark accuracy; \gls{cpu} throughput & No \\
Wang et al.~\cite{slmquant} & Evaluation & Sub-billion models, Qwen2.5-0.5B among them & Methods tuned on large models, applied to small ones & Generic benchmark accuracy & No \\
\midrule
This work & Evaluation & Three fine-tuned models of 0.36--1.2~B & Q8\_0 and Q4\_K\_M against fp16 & \gls{em}, intent and slot F1, schema validity, false-command rate, on one task; throughput on the board & Yes \\
\bottomrule
\end{tabularx}
\end{table}

\paragraph{Critical comparison.} Read down its columns, Table~\ref{tab:soa-quantisation}
exhibits four patterns. On bit assignment, every method uses data to decide where precision goes,
whether calibration inputs, activation statistics or the weights' own outliers, whereas the runtime
deployed here assigns bits by a fixed rule; the deployed format is therefore not an instance of any
method in the table, and no method's reported cost transfers to it. On the cost metric, the methods
report perplexity, with GPTQ adding zero-shot accuracy on generic suites, while the evaluations
report task accuracy; the shift from a distributional measure to a task measure is recent and is
made on generic benchmarks only. On scale, the evaluations that score task accuracy work at eight
billion parameters and above, with one reaching down to 1.5~billion, and the one study of
sub-billion models finds that they behave differently under quantisation, which is a reason not to
extrapolate downwards from the others. On decoding, no study in the table quantises a model and
then scores it under a grammar; every accuracy figure is obtained with the decoder unconstrained.
None of these studies scores a structured command-parsing task by \gls{em}, intent and slot F1, or
schema validity, and none decodes under a grammar. Chapter~\ref{chap:results} measures the
accuracy cost of the whole deployment pipeline on that task, for the three quantised models at
both levels, as part of Contribution~C3. Section~\ref{sec:quantisation-procedure} explains why
that pipeline cost is not the cost of quantisation alone, and why it is not to be read against the
studies above.

\section{Constrained decoding}
\label{sec:soa-constrained-decoding}
Grammar-constrained generation, the token-masking mechanism of Section~\ref{sec:bg-gbnf}, has
been applied to structured language-processing tasks without fine-tuning~\cite{geng2023}. The
literature around it is of three kinds: the constructions, which make the mask correct and cheap;
the runtime implementations, which apply it; and the cost studies, which ask what constraining
the decoder does to the model's output. The constructions compile a regular expression or a
context-free grammar into an automaton over the vocabulary and mask the decoder against
it~\cite{willard2023,koo2024}, with correctness results for the language classes such a
construction covers~\cite{koo2024}. \Gls{gbnf} is the formalism \texttt{llama.cpp} implements it
with~\cite{llamacpp}. The mechanism differs from prompting a model to produce well-formed output
and then checking the result: the grammar acts inside the sampling loop, before an invalid token
can be emitted, not after a full sequence has been generated.

That difference in mechanism is also a difference in the kind of claim that can be made about
structural validity. Without a grammar, whether a model's output parses against a schema is an
empirical question, answered by counting how often it does so on some test set, and the answer
can change whenever the input distribution changes. With a grammar written from the same schema,
the question is closed by construction: no string outside the grammar's language is reachable, on
any input, provided the generation cap exceeds the longest string that language contains. The
bounded grammar of Section~\ref{sec:schema} makes that length finite. Constrained decoding thus
moves structural validity from a rate that fine-tuning and a favourable test set can inflate to a
property the decoder enforces independently of both. The grammar ablation of
Chapter~\ref{chap:results} decodes the same fine-tuned models with and without the grammar, and so
measures how often they leave the format when nothing prevents them.

That property has a cost, which is why the ablation reports accuracy as well as validity.
Constraining the decoder alters the distribution the model samples from: outputs remain
grammatical, but their relative likelihoods no longer track what the unconstrained model would
have assigned, so a grammar can make a structurally valid command more likely without making the
correct one more likely~\cite{park2024}. Format restriction has separately been observed to
degrade accuracy on reasoning tasks, more so as the restriction tightens, while classification
tasks were in some cases helped by it~\cite{tam2024}. Table~\ref{tab:soa-constrained} arranges the
six works by what they contribute and what they measure.

\begin{table}[htbp]
\centering
\footnotesize
\caption[Constrained-decoding constructions, implementations and cost studies]{Constrained-decoding
constructions, implementations and cost studies. \emph{Fine-tuned to the format}: whether the
models evaluated had been trained to produce the constrained format before the constraint was
applied. A dash marks a property this review did not establish from the source. The last row is
this work.}
\label{tab:soa-constrained}
\begin{tabularx}{\textwidth}{@{}>{\raggedright\arraybackslash}p{2.3cm} >{\raggedright\arraybackslash}p{2.5cm} >{\raggedright\arraybackslash}X >{\raggedright\arraybackslash}p{1.8cm} >{\raggedright\arraybackslash}X@{}}
\toprule
Work & Kind & Constraint and mechanism & Fine-tuned to the format & Effect on accuracy measured \\
\midrule
Geng et al.~\cite{geng2023} & Application & Context-free grammar; token masking on structured NLP tasks & No & --- \\
Willard and Louf~\cite{willard2023} & Construction & Regular expression or context-free grammar compiled to a finite-state index over the vocabulary & --- & --- \\
Koo et al.~\cite{koo2024} & Construction & Automata for regular and deterministic context-free languages; provably correct masks & --- & --- \\
Park et al.~\cite{park2024} & Cost study & Grammar-constrained decoding, shown to distort the model's distribution; an aligned alternative proposed & --- & Yes, as divergence from the unconstrained distribution \\
Tam et al.~\cite{tam2024} & Cost study & Format restriction of increasing strictness & No & Yes: reasoning accuracy falls as restriction tightens; classification sometimes helped \\
\texttt{llama.cpp} \gls{gbnf}~\cite{llamacpp} & Implementation & \gls{bnf} dialect parsed per request and applied per step; cost depends on rule shape & --- & --- \\
\midrule
This work & Ablation & Bounded \gls{gbnf} grammar written from the command schema; decoded with and without it & Yes & Yes: \gls{em}, F1 and schema validity with and without the grammar \\
\bottomrule
\end{tabularx}
\end{table}

\paragraph{Critical comparison.} Table~\ref{tab:soa-constrained} separates the works along three
axes. On what they contribute, the constructions establish that a mask can be built correctly and
applied cheaply, the application demonstrates that it works without fine-tuning, and the cost studies show
that it is not free; the three kinds do not cite one another's measurements, so the literature
establishes correctness and cost separately, on different models and tasks. On fine-tuning, no work
in the table evaluates a model that had already been trained to produce the constrained format:
the application and the format-restriction study use models prompted into the format, and the
constructions report no task evaluation at all. What a grammar adds once fine-tuning has taught a
model the format is therefore not reported anywhere in the table. On the direction of the effect,
the two cost studies disagree in kind: the distortion result establishes that the constrained
distribution differs from the unconstrained one without stating how much accuracy that costs on
any task, and the format-restriction study finds the sign of the effect task-dependent. The size
of either effect on a given task does not follow from that literature. Of the two task types the
format-restriction study compares, classification is the closer to parsing a command into one of
ten intents. Chapter~\ref{chap:results} therefore reports accuracy with and without the grammar,
not structural validity alone, on models that were fine-tuned to the format first.

\section{Spoken-language understanding for robotics}
\label{sec:soa-slu}
Parsing an utterance into an intent and a set of slot values is an established formulation in
spoken-language understanding, not one this thesis introduces. The literature on it is of two
kinds: surveys that fix the formulation, and corpora with their baselines that instantiate it. A
survey of the field treats the extraction of this semantic frame as the object of the task, and
classifies methods by whether they model intent and slots separately or jointly~\cite{qin2021}.
MASSIVE, a dataset of one million utterances in 51 languages labelled with 60 intents and 55 slot
types, is built around the same formulation~\cite{massive}. Its baselines, encoders of 258 to
580~million parameters, reach locale-averaged intent accuracies of 85.1--86.1\% and slot F1 of
73.6--76.8\%, but \gls{em} of only 63.7--66.6\%~\cite{massive}. For the single-utterance commands
considered here, a spoken instruction to a robot has the same shape: one intent from a fixed set
and a handful of slot values. Language-model control of robots and aircraft more broadly is
reviewed in the \emph{M\'emoire d'Ing\'enieur}. Table~\ref{tab:soa-slu} sets the survey, the
corpus and this work on one set of axes.

\begin{table}[htbp]
\centering
\footnotesize
\caption[Intent-and-slot parsing: formulation, benchmark and this work]{Intent-and-slot parsing
as formulated by the survey, instantiated by the benchmark, and deployed here. A dash marks a
property this review did not establish from the source. The last row is this work.}
\label{tab:soa-slu}
\begin{tabularx}{\textwidth}{@{}>{\raggedright\arraybackslash}p{2.4cm} >{\raggedright\arraybackslash}p{1.6cm} >{\raggedright\arraybackslash}X >{\raggedright\arraybackslash}X >{\raggedright\arraybackslash}p{2.4cm} >{\raggedright\arraybackslash}p{1.7cm}@{}}
\toprule
Work & Kind & Formulation and models & Metrics reported & Hardware and latency budget & Output constraint \\
\midrule
Qin et al.~\cite{qin2021} & Survey & Semantic frame of intent and slots; methods sorted into single and joint models & --- & --- & --- \\
MASSIVE~\cite{massive} & Corpus and baselines & 60 intents, 55 slot types, 51 languages; encoder baselines of 258--580~M parameters & Intent accuracy 85.1--86.1\%; slot F1 73.6--76.8\%; \gls{em} 63.7--66.6\% & Data-centre \glspl{gpu}; no budget & None \\
\midrule
This work & Corpus, models and deployment & Ten intents and a bounded slot vocabulary for a five-drone swarm; decoder-only models of 0.36--1.2~B, fine-tuned & \gls{em}, intent and slot F1, schema validity, false-command and safe-failure rates & Raspberry~Pi~5 \gls{cpu}; 1{,}100~ms decode allowance & Bounded \gls{gbnf} grammar \\
\bottomrule
\end{tabularx}
\end{table}

\paragraph{Critical comparison.} What differs across Table~\ref{tab:soa-slu} is the setting, not
the shape of the task. On formulation, the survey and the corpus agree: one intent from a closed
set and a set of typed slots, with the joint model as the modern default. On the models, the
benchmark's baselines are encoders that classify and label, whereas this work uses a decoder that
generates the frame as text, which is what makes a grammar applicable to it. On the metric, the
benchmark's lowest figure is \gls{em}, and it is the metric a controller depends on, since a
command with one wrong slot is a wrong command; the benchmark reports it but does not optimise for
it, and its baselines are trained on data-centre \glspl{gpu}, decoded without a grammar, and
scored with no latency or memory limit~\cite{massive}. On the deployment, this work's model has to
run offline, at sub-billion-parameter scale, within the throughput budget of a Raspberry~Pi~5, and
its output is acted on by the swarm controller after automatic validation, with no person reading
it first. No existing corpus covers this command vocabulary (Section~\ref{sec:dataset}). That the
task shape is well studied does not show that it can be solved at this size and within this
budget, with an error rate low enough for output that no person reads. The evaluation protocol of
Chapter~\ref{chap:method} is built to answer that question.

\section{Research gaps}
\label{sec:research-gaps}

\paragraph{Synthesis.} Each of the four literatures reviewed in this chapter answers its own
question. Benchmarking on \glspl{sbc} has found that such boards reliably serve models of up to
about 1.5~billion parameters, and that on one board the choice of runtime changed throughput by up
to four times~\cite{sbc2025}. Post-training quantisation is the established way to shrink a trained
model~\cite{gptq}, with on-device deployment among its motivations~\cite{awq}. Its cost is reported as perplexity and generic zero-shot accuracy by
the papers introducing the methods~\cite{spqr,gptq}, and as accuracy on general-purpose suites by
the studies evaluating them, including at sub-billion scale~\cite{kurtic2025,kurt2026,slmquant}.
Grammar-constrained decoding makes a model's output structurally valid by construction rather than
by the accuracy of the model that produced it~\cite{willard2023,koo2024}, at costs its own
literature documents: a distorted sampling distribution~\cite{park2024} and, on reasoning tasks,
lower accuracy~\cite{tam2024}. Intent and slot parsing
is an established task formulation, with a large multilingual benchmark built for
it~\cite{massive}. Some pairs of these literatures have already been combined: quantisation with
\gls{cpu} inference and downstream accuracy~\cite{kurt2026}, quantisation with small
models~\cite{slmquant}, and format restriction with task accuracy~\cite{tam2024}. No study
reviewed here combines all four. The four critical comparisons above leave three gaps, stated
here in the order the contributions of Chapter~\ref{chap:introduction} close them.

\paragraph{Gap G1, structural validity after fine-tuning.} The constrained-decoding studies
reviewed here (Table~\ref{tab:soa-constrained}) evaluate grammars and format restrictions on
models that were not fine-tuned to the target format, so what a grammar still adds once
fine-tuning has taught a model that format is not reported, and neither is what it costs in
accuracy on a task of the classification kind. Contribution C1, the frozen command schema and its
matched, bounded \gls{gbnf} grammar, makes structural validity on this task a property of the
decoder. The grammar ablation, run with the harness of Contribution C3, decodes the same
fine-tuned, quantised models with and without that grammar, and so reports what the grammar adds
once fine-tuning has taught the models the format.

\paragraph{Gap G2, the command corpus.} Intent-and-slot benchmarks such as
MASSIVE (Table~\ref{tab:soa-slu}) are scored on data-centre hardware without a latency budget or a
grammar, and no existing corpus covers this command vocabulary. Contribution C2, the label-first
dataset (each label written before the sentences that express it), supplies what a general-purpose
corpus does not: training and test data for this command vocabulary, split at the template-family
level, with a recorded golden set on which the deployed configuration is scored.

\paragraph{Gap G3, deployment cost on a task.} The single-board-computer
benchmark (Table~\ref{tab:soa-edge}) reports throughput without reference to any task, so it
cannot say what a configuration that meets a throughput floor costs in task accuracy. The
quantisation studies (Table~\ref{tab:soa-quantisation}) report accuracy on general-purpose
benchmarks with the decoder unconstrained, not \gls{em}, intent and slot F1, or schema validity on
a command-parsing task, and those that score tasks at all work at eight billion parameters and
above. Choosing a model to deploy on this hardware requires all four quantities measured together,
on the configuration that is actually deployed. Contribution C3, the quantised-inference harness,
measures \gls{em}, intent and slot F1, schema validity, and false-command rate for three of the
four fine-tuned models at Q8\_0 and Q4\_K\_M, with accuracy scored under the deployed runtime and
grammar and throughput timed on the Raspberry~Pi~5 at Q4\_K\_M. It reports the accuracy cost of
the whole deployment pipeline on this task at sub-billion scale, next to the throughput that
single-board-computer benchmarking measures alone.

Together the three contributions supply the measurements from which Chapter~\ref{chap:results}
answers RQ1 for this task and this device. Chapter~\ref{chap:method} specifies how each
contribution was built, and Chapter~\ref{chap:results} reports what it measured.
